\section*{Erd\H{o}s Problem 641: cycles on the same vertex set}

\paragraph{Statement.}
Does there exist, for every positive integer \(k\), a number \(f(k)\)
such that every graph of chromatic number at least \(f(k)\) contains
\(k\) pairwise edge-disjoint cycles on exactly the same vertex set?

\paragraph{Claim and external input.}
The answer is negative for every \(k\geq2\), and \(f(1)=3\).
The decisive construction is external: Janzer, Steiner, and Sudakov,
\emph{Chromatic number and regular subgraphs},
\url{https://arxiv.org/abs/2410.02437}, Theorem 1.2, prove that for
some absolute \(c>0\) and all sufficiently large \(n\) there is an
\(n\)-vertex graph \(G_n\) with no \(4\)-regular subgraph and
\[
 \chi_f(G_n)\geq c\,\frac{\log\log n}{\log\log\log n}.
\]
We use that construction as a stated theorem, not as a construction
proved anew in this attempt. The following reduction checks all the
quantifiers needed for \url{https://www.erdosproblems.com/641}.

\paragraph{Reduction.}
An ordinary proper \(r\)-colouring partitions the vertices into \(r\)
independent sets. Giving each of these sets weight \(1\) is a feasible
fractional colouring of total weight \(r\). Hence
\(\chi_f(G)\leq\chi(G)\).
Since \(\log\log n/\log\log\log n\) tends to infinity, the graphs
in the theorem have unbounded ordinary chromatic number.

If two edge-disjoint cycles have the same vertex set \(S\), every
vertex of their union has two incident edges from each cycle.
The four edges are distinct, so their union is a \(4\)-regular
subgraph on \(S\). Additional ambient edges are irrelevant: the
subgraph is not required to be induced. Thus none of the graphs
\(G_n\) contains such a pair.
For any proposed finite \(f(k)\), \(k\geq2\), choose \(n\) large
enough that \(\chi_f(G_n)\geq f(k)\). A family of \(k\) cycles would
contain a forbidden pair. This disproves the proposed threshold.

For \(k=1\), a graph with no cycle is a forest and is two-colourable
by the parity of distance from a root in each component. Thus
\(\chi(G)\geq3\) forces a cycle. A single edge shows that a threshold
at most \(2\) would fail.

\paragraph{Why a minimum-degree argument does not settle it.}
If \(\chi(G)\geq r+1\), repeatedly removing vertices of degree at most
\(r-1\) cannot remove the whole graph: otherwise colouring in reverse
order would use at most \(r\) colours. Therefore such a graph has a
subgraph of minimum degree at least \(r\).
This observation does not produce a regular subgraph. The external
construction shows precisely that even arbitrarily large chromatic
number cannot force the required \(4\)-regular structure.

\paragraph{Scope.}
The disproof is complete relative to the explicitly stated
Janzer--Steiner--Sudakov construction. Its probabilistic construction
is not re-proved here.

\paragraph{Completion estimate.}
\textbf{COMPLETION: 100\%}
